\chapter{La plateforme multigérants}\label{ch:fm:the-multi-manager-platform}

La brochure réglementaire d'une grande plateforme multigérants énumère, parmi les charges supportées par ses fonds, les salaires, les primes et la
rémunération liée à la performance des gérants de portefeuille qui dirigent ses équipes, les honoraires de leurs recruteurs, leurs ordinateurs et leurs bureaux. Elle appelle ce dispositif «~
refacturation des charges~». Elle ajoute que le gérant de portefeuille touche une part du profit de sa propre équipe, «~sans tenir compte de la
performance des autres gérants de portefeuille ou des Millennium Partners Funds en général~», et que les gérants dont la performance est positive
«~toucheront une rémunération liée à la performance même si les rendements globaux des Millennium Partners Funds sont négatifs~». L'investisseur verse la
prime de chaque équipe gagnante, y compris les années où le fonds perd de l'argent. Ce chapitre chiffre cette promesse.

\section{Pods, le centre et la promesse de la plateforme}

Une plateforme multigérants (livre~1, chapitre~1) est un hedge fund géré en autant d'équipes indépendantes, chacune étant un pod (livre~8, chapitre~28),
sous l'égide d'un centre qui alloue le capital, fixe les limites de risque, compense les expositions et couvre leur somme, et paie l'infrastructure. Sa promesse aux investisseurs est
une diversification des talents~: cinquante équipes aux ratios de Sharpe modestes et faiblement corrélées forment un fonds dont le ratio est élevé. Sa promesse aux gérants de portefeuille
est une part, fixée par une formule, de ce que gagne leur propre équipe, versée quoi que fassent les autres.

\begin{proposition}[Diversification des talents]\label{prop:fm:the-multi-manager-platform:div}
Soient $n$ pods de même volatilité $\sigma$, de même ratio de Sharpe $S$ et de même corrélation deux à deux $\rho$, et supposons que le fonds les détienne à poids égaux.
Le ratio de Sharpe du fonds est
\[
S_{\mathrm{fund}}=S\sqrt{\frac{n}{1+(n-1)\rho}}\;\longrightarrow\;\frac{S}{\sqrt\rho}\quad(n\to\infty).
\]
\end{proposition}

\begin{proof}
L'espérance de rendement du fonds vaut $S\sigma$ et sa variance $\sigma^2(1+(n-1)\rho)/n$ par unité de capital, répartie également~; il suffit de diviser. La limite s'obtient
lorsque le terme $(n-1)\rho$ domine.
\end{proof}

\begin{center}
\small
\begin{tabular}{@{}lrrrr@{}}
\toprule
ratio de Sharpe du fonds (pods à $S=1$) & $n=10$ & $n=25$ & $n=50$ & $n=100$\\
\midrule
$\rho=0.05$ & 2.63 & 3.37 & 3.81 & 4.10\\
$\rho=0.10$ & 2.29 & 2.71 & 2.91 & 3.03\\
$\rho=0.20$ & 1.89 & 2.08 & 2.15 & 2.19\\
\bottomrule
\end{tabular}
\end{center}

Le tableau est à la fois l'argumentaire économique de la plateforme et sa limite. La corrélation, et non le nombre d'équipes, fixe le plafond~: à $\rho=0.1$, il est
$1/\sqrt{0.1}=3.16$ fois le ratio de Sharpe d'un pod, et cinquante pods atteignent déjà 2.91 de ce plafond. Au-delà de quelques dizaines d'équipes, ajouter des pods rapporte peu, sauf
si leur corrélation avec celles déjà en place est plus faible~; réduire la corrélation, en recrutant des stratégies différentes ou en couvrant leurs expositions communes
(\cref{def:fm:the-multi-manager-platform:centre}), rapporte davantage que le recrutement.

\begin{omfigure}
\begin{tikzpicture}[font=\footnotesize,>=stealth,box/.style={draw,rounded corners,minimum height=0.62cm,align=center,fill=omDef!10}]
\node[box,minimum width=2.4cm] (inv) at (0,3.6) {investisseurs};
\node[box,minimum width=3.0cm] (fund) at (0,2.2) {le fonds\\{\scriptsize capital, coûts, frais}};
\node[box,fill=omProp!12,minimum width=3.0cm] (centre) at (5.2,2.2) {le centre\\{\scriptsize allocation, risque, book central}};
\foreach \x/\n in {-3.6/1,-1.8/2,0/3,1.8/4,3.6/5} {\node[box,fill=gray!12,minimum width=1.5cm] (p\n) at (\x,0.4) {pod \n};}
\node[font=\scriptsize] at (5.2,0.4) {\dots\ cinquante pods};
\draw[->] (inv) -- node[right,font=\scriptsize] {capital} (fund);
\foreach \n in {1,2,3,4,5} {\draw[->] (fund.south) -- (p\n.north);}
\draw[<->,dashed] (fund.east) -- (centre.west);
\draw[->,omThm] (p1.south) -- ++(0,-0.5) -| node[pos=0.25,below,font=\scriptsize] {chaque équipe touche une part de son propre profit} (p5.south);
\node[font=\scriptsize,omThm,align=left,anchor=west] at (-6.0,2.2) {pass-through~:\\coûts de toutes les équipes\\et rémunérations};
\end{tikzpicture}

\omcaption{Une plateforme multigérants. Le capital des investisseurs est alloué aux pods par le centre, qui couvre ou renforce leurs expositions nettes
par son propre book~; le fonds supporte l'ensemble des coûts et des rémunérations de chaque équipe (pass-through), et chaque équipe est rémunérée sur son propre profit. Modèle~:
\texttt{firm.podshop}.}\label{fig:fm:the-multi-manager-platform:structure}
\end{omfigure}

\begin{definition}[Frais refacturés]\label{def:fm:the-multi-manager-platform:pass}
Les \emph{frais refacturés}\index{frais refacturés} (pass-through fee) sont la somme que le fonds facture à ses investisseurs au titre des coûts réels
de son fonctionnement, à savoir la rémunération de ses équipes d'investissement et de son personnel, de ses données, de ses systèmes et de ses bureaux, au lieu
ou en complément de frais de gestion fixes, sur lesquels le gérant prélèverait leur rémunération.
\end{definition}

\begin{definition}[Taux de reversement]\label{def:fm:the-multi-manager-platform:payout}
Le \emph{taux de reversement}\index{taux de reversement} (payout rate) d'une équipe est la part de son propre résultat de trading, nette de ses coûts et des pertes reportées, qui
est reversée à l'équipe sous forme de rémunération liée à la performance.
\end{definition}

\begin{dated}{2026-09}{Ce que dit la brochure d'une plateforme}\label{dat:fm:the-multi-manager-platform:adv}
Millennium Management LLC, Form ADV Part 2A (25 août 2026)~: ses fonds «~supportent en général, directement ou indirectement, toutes les charges engagées
dans le cadre de~» leur fonctionnement, un dispositif qu'elle nomme «~refacturation des charges~», y compris la rémunération des gérants de portefeuille et des autres
salariés, honoraires des recruteurs, équipements et bureaux~; ce dispositif est «~distinct de, et en plus de~» la rémunération liée à la performance et aux encours que perçoit le
conseiller, et elle est déduite avant celle-ci. La rémunération d'un gérant de portefeuille est «~généralement déterminée comme un pourcentage des profits réalisés par
ce gérant de portefeuille au cours de l'année civile précédente~»~; les pertes sont reportées~; il n'y a «~généralement aucun report en arrière ni récupération~». Le
conseiller alloue et réalloue le capital entre les gérants de portefeuille et effectue des investissements directs avec le capital des fonds, y compris des «~couvertures, ou
opérations contraires visant à réduire certaines expositions~» et des transactions qui accroissent l'exposition «~aux positions nettes détenues par un nombre
de gérants de portefeuille~».
\end{dated}

\section{Frais refacturés et facture de l'investisseur}

La plateforme du chapitre (à titre illustratif) compte 50 pods, chacun avec un ratio de Sharpe de 1.0 et une volatilité de 20~\% par an sur sa part du
capital du fonds, corrélés à 0.1 par un facteur commun. Les équipes touchent 20~\% de leur propre profit, pertes reportées~; une équipe qui perd plus
que 10~\% de son capital en un an est remplacée par une nouvelle. Les coûts des pods, 3~\% par an, et le 1~\% de la plateforme sont refacturés~; le gérant prélève
une allocation de performance de 20~\% sur ce qui reste. Sur vingt années simulées et vingt graines aléatoires, cent dollars de capital rapportent chaque année
(\cref{fig:fm:the-multi-manager-platform:waterfall})~: 20.09~\$ brut, dont 4.27~\$ reviennent aux équipes, 4.00~\$ aux coûts et 2.38~\$ au gérant,
laissant à l'investisseur 9.44~\$. Le fonds a perdu de l'argent dans 9 de ses 400 années.

\begin{omfigure}
\begin{tikzpicture}
\begin{axis}[width=10.6cm,height=5.2cm,ybar stacked,bar width=22pt,ymin=0,ymax=22,ylabel={\small \$ par 100~\$ de capital par an},
  xtick={0,1,2,3,4},xticklabels={brut,{reversements aux équipes},{coûts refacturés},gérant,investisseur},
  xticklabel style={font=\scriptsize,text width=1.75cm,align=center},tick label style={font=\footnotesize},
  table/col sep=comma,enlarge x limits=0.12,grid=major,grid style={gray!20},axis x line*=bottom]
\addplot[draw=none,fill=none] table[x=k,y=bottom]{figdata/desk/03-the-multi-manager-platform/waterfall.csv};
\addplot[fill=omThm!50,draw=omThm] table[x=k,y=cost]{figdata/desk/03-the-multi-manager-platform/waterfall.csv};
\addplot[fill=omDef!55,draw=omDef] table[x=k,y=total]{figdata/desk/03-the-multi-manager-platform/waterfall.csv};
\end{axis}
\end{tikzpicture}

\omcaption{Où va le profit brut d'une plateforme, par 100~\$ de capital et par an~: reversements aux équipes sur le profit propre de chaque pod, coûts refacturés,
l'allocation de performance du gérant et ce que garde l'investisseur. Cinquante pods de ratio de Sharpe 1.0, vingt années, vingt graines~; conditions illustratives.
Données~: \texttt{fm\_platform.waterfall}.}\label{fig:fm:the-multi-manager-platform:waterfall}
\end{omfigure}

Un fonds classique qui prélève 2~\% du capital et 20~\% du profit sur les mêmes pods laisserait à ses investisseurs 14.47~\$, en leur prélevant 5.62~\$. Il ne pourrait pas
exister~: sur ces 5.62~\$ de frais, son gérant devrait verser 4.27~\$ aux équipes et 4.00~\$ de coûts, soit une perte de 2.65~\$ pour chaque centaine de dollars. Le
pass-through est la structure de frais dont une plateforme a besoin dès que les talents et l'infrastructure coûtent plus qu'un forfait fixe ne peut couvrir~; la question de l'investisseur est
de savoir si ce qui reste après ces coûts, soit 9.44~\$ ici, vaut le risque du fonds.

\omcode{code/firm/podshop/firm_podshop.py}{48}{66}{Le versement de chaque équipe sur son propre profit, les pertes reportées et les équipes remplacées~; le
coût de compensation est ce que coûtent les versements au-delà du même taux appliqué au total du fonds.}\label{lst:fm:the-multi-manager-platform:payouts}

\section{Rémunérer les gagnants~: le problème de la compensation}

\begin{definition}[Risque de compensation des résultats]\label{def:fm:the-multi-manager-platform:netting}
Le \emph{risque de compensation des résultats}\index{risque de compensation des résultats} (netting risk) d'un fonds multigérants est le coût pour ses investisseurs de verser à chaque équipe une part de son propre profit
sans le compenser avec les pertes des autres équipes~: l'écart entre la somme des versements aux équipes et le même taux de versement appliqué au
profit total du fonds.
\end{definition}

\begin{proposition}[Le prix de la rémunération des gagnants]\label{prop:fm:the-multi-manager-platform:netting}
Soit le profit annuel d'un pod $X\sim\mathcal N(S\sigma,\sigma^2)$, de ratio de Sharpe $S$, et l'on verse un taux $p$ sur $\max(X,0)$ chaque année, sans
report des pertes. Si le profit total du fonds est positif, le coût de compensation par pod vaut
\[
p\,\bigl(\E[\max(X,0)]-\E[X]\bigr)=p\,\sigma\,\bigl(\varphi(S)-S\,\Phi(-S)\bigr),
\]
soit une part $p\,(\varphi(S)-S\Phi(-S))/S$ du profit espéré du pod.
\end{proposition}

\begin{proof}
$\E[\max(X,0)]=\mu\Phi(\mu/\sigma)+\sigma\varphi(\mu/\sigma)$ pour $X\sim\mathcal N(\mu,\sigma^2)$~; on retranche $\mu=S\sigma$ et l'on utilise
$1-\Phi(S)=\Phi(-S)$. Lorsque le total du fonds est positif, le versement sur le total vaut $p$ fois la somme des $X$, dont l'espérance vaut $p$ fois
la somme des $\E[X]$.
\end{proof}

À $p=20~\%$, le coût de compensation représente 1.7~\% du profit brut pour des pods de ratio de Sharpe 1, 7.9~\% à 0.5 et 22.9~\% à 0.25
(\cref{fig:fm:the-multi-manager-platform:netting}). Reporter les pertes supprime l'essentiel de ce coût pour une équipe qui reste~: dans la plateforme du chapitre,
où aucune équipe n'est jamais remplacée, le coût de compensation vaut 0.1~\% du profit brut. Remplacer les équipes perdantes le rétablit, car la perte reportée d'une équipe remplacée disparaît
avec elle et son successeur part de zéro~: lorsque les équipes sont remplacées après une perte de 10~\%, le coût simulé vaut 1.3~\% du profit brut à un ratio de Sharpe de
1 et 6.2~\% à 0.5, proche de la borne de la formule. Le risque de compensation d'une plateforme dépend donc de la qualité de ses équipes et du rythme auquel elle
les renouvelle, et non du seul taux de versement.

\begin{omfigure}
\begin{tikzpicture}
\begin{axis}[width=10.5cm,height=5.2cm,xlabel={\small ratio de Sharpe de chaque pod},ylabel={\small coût de compensation (\% du profit brut)},
  xmin=0.2,xmax=2.05,ymin=0,ymax=25,tick label style={font=\footnotesize},grid=major,grid style={gray!20},table/col sep=comma,
  legend style={font=\scriptsize,at={(0.5,-0.3)},anchor=north,/tikz/every even column/.append style={column sep=8pt}},legend columns=2]
\addplot[omThm,thick,mark=*] table[x=sr,y=closed]{figdata/desk/03-the-multi-manager-platform/netting.csv};
\addplot[omDef,thick,dashed,mark=square*] table[x=sr,y=simulated]{figdata/desk/03-the-multi-manager-platform/netting.csv};
\legend{sans report des pertes (\cref{prop:fm:the-multi-manager-platform:netting}),report des pertes et équipes remplacées après une perte de 10~\%}
\end{axis}
\end{tikzpicture}

\omcaption{Le coût de rémunération de chaque équipe sur son propre profit, au-delà du même taux de 20~\% appliqué au total du fonds, en part du profit brut, en fonction du
ratio de Sharpe des pods~: la forme fermée sans report des pertes, et la plateforme du chapitre, cinquante pods corrélés à 0.1, vingt années, vingt graines.
Données~: \texttt{fm\_platform.netting\_vs\_sr}.}\label{fig:fm:the-multi-manager-platform:netting}
\end{omfigure}

\begin{example}[Deux pods]\label{ex:fm:the-multi-manager-platform:two}
Deux pods de 100 millions de dollars chacun~: l'un gagne 10 millions de dollars, l'autre perd 10 millions de dollars. Le profit du fonds est nul. À un taux de versement de 20~\%, la première
équipe reçoit 2 millions de dollars et la seconde rien~: les investisseurs versent 2 millions de dollars, auxquels s'ajoutent les coûts des deux pods, pour une année où le fonds n'a rien gagné.
Si l'équipe perdante reste, ses 10 premiers millions de dollars de profit ultérieur ne lui rapportent rien~; si elle est remplacée, son successeur est rémunéré dès le premier dollar.
\end{example}

\section{Allocation de capital et règles de drawdown}

Le centre alloue du capital aux équipes et le leur retire. Le livre~8, chapitre~28, a simulé l'arbitrage central~: une limite de drawdown
(définie là-bas) arrête les bonnes équipes en mauvaise passe et laisse courir les mauvaises, et son taux de faux arrêts est fixé par la limite exprimée en unités de la volatilité propre
de l'équipe. Les plateformes décrivent leurs règles en termes généraux~; les seuils précis des sociétés nommées ne figurent pas dans leurs documents publics, et ce
livre ne reprend pas de chiffres de seconde main. Le modèle utilise donc une échelle illustrative à deux paliers~: le capital d'une équipe est réduit de moitié après un drawdown de 10~\%
dans l'année et retiré à 15~\%, puis restitué au début de l'année suivante.

\begin{method}[Évaluer l'échelle de drawdown d'une plateforme]\label{met:fm:the-multi-manager-platform:ladder}
\begin{enumerate}
\item Exprimer chaque seuil en unités de la volatilité de l'équipe sur son capital (une réduction de 10~\% équivaut à la moitié de la volatilité annuelle à 20~\%).
\item Simuler les pods de la plateforme avec et sans l'échelle, sur les mêmes tirages aléatoires.
\item Indiquer la part des années-équipes réduites et arrêtées, ainsi que le rendement, la volatilité et le ratio de Sharpe du fonds dans les deux cas.
\item Indiquer séparément ce que l'échelle fait aux équipes perdantes (chapitre~8)~: une échelle qui n'améliore rien dans le monde des bonnes équipes doit se justifier par son coût
dans celui des mauvaises équipes.
\end{enumerate}
\end{method}

Dans la plateforme du chapitre, l'échelle réduit 88.4~\% des années-équipes et en arrête 26.6~\%~: avec une volatilité de 20~\%, un drawdown de 10~\% dans l'année est
ordinaire. Le rendement du fonds passe de 20.1~\% à 13.4~\% et sa volatilité de 6.9~\% à 4.8~\%~; son ratio de Sharpe baisse de 2.92 à 2.77. Chaque
pod de ce modèle possède un edge réel, si bien que l'échelle ne peut que coûter~; le chapitre~8 mesure ce qu'elle rapporte lorsque certains edges disparaissent.

\section{Book central et surcouches}

\begin{definition}[Book central, surcouche factorielle]\label{def:fm:the-multi-manager-platform:centre}
Le \emph{book central}\index{book central} (centre book) d'une plateforme est le book que le centre traite lui-même avec le capital du fonds, aux côtés des pods~: pour couvrir
les expositions que partagent les pods, pour renforcer des positions que détiennent beaucoup de pods, ou pour traiter ses propres stratégies. Une \emph{surcouche factorielle}\index{surcouche factorielle} (factor overlay) est une
position du book central qui compense la somme des expositions des pods à un facteur commun, de sorte que le fonds conserve les rendements spécifiques des pods sans
leur risque commun.
\end{definition}

Cinquante pods corrélés à 0.1 donnent un fonds dont la volatilité est dominée par ce qu'ils partagent~: $\sigma\sqrt{(1+49\rho)/50}$ vaut 6.9~\% pour des pods de 20~\%
de volatilité, contre 2.7~\% si la part commune est couverte. Dans le modèle, la surcouche laisse le rendement moyen du fonds à 20.1~\% et fait passer son ratio de Sharpe
de 2.9 à 7.5 (\cref{fig:fm:the-multi-manager-platform:overlay}). C'est une borne supérieure~: ici, tout le risque commun tient en un seul facteur que le
centre peut traiter. Les pods réels partagent des positions encombrées et de la liquidité en plus des facteurs de marché et de style (le livre~7, chapitre~28), et seule
la part négociable peut être couverte~; la brochure du hook décrit les deux usages des transactions propres du centre~: couvrir les expositions compensées et renforcer
celles-ci.

\begin{omfigure}
\begin{tikzpicture}
\begin{axis}[width=10.5cm,height=5.2cm,xlabel={\small année},ylabel={\small rendement cumulé (\%)},xmin=0,xmax=20,ymin=-20,ymax=420,
  tick label style={font=\footnotesize},grid=major,grid style={gray!20},table/col sep=comma,
  legend style={font=\scriptsize,at={(0.5,-0.3)},anchor=north,/tikz/every even column/.append style={column sep=8pt}},legend columns=2]
\addplot[omThm,thick] table[x=year,y=base]{figdata/desk/03-the-multi-manager-platform/overlay.csv};
\addplot[omDef,thick,dashed] table[x=year,y=hedged]{figdata/desk/03-the-multi-manager-platform/overlay.csv};
\legend{pods seuls,avec la surcouche factorielle}
\end{axis}
\end{tikzpicture}

\omcaption{La plateforme du chapitre sur vingt ans (graine 1), brute de frais, avec et sans une surcouche du \omterm{def:fm:the-multi-manager-platform:centre}{book central} qui couvre l'exposition sommée des pods
à leur facteur commun. La surcouche conserve la moyenne et supprime les oscillations communes. Données~: \texttt{fm\_platform.sim}, \texttt{firm.podshop.overlay}.}\label{fig:fm:the-multi-manager-platform:overlay}
\end{omfigure}

\begin{remark}[À qui appartient le P\&L du book central]\label{rem:fm:the-multi-manager-platform:owner}
Le profit d'un \omterm{def:fm:the-multi-manager-platform:centre}{book central} appartient au fonds, et non à une équipe, si bien qu'il ne donne lieu à aucun reversement d'équipe~; ses coûts sont refacturés comme les autres. Lorsque le centre
ajoute des positions que de nombreux pods détiennent, il accroît la concentration que les pods créent déjà, et son risque doit figurer dans les mêmes limites de risque que les leurs
(chapitre~7).
\end{remark}

\section{Tutoriel~: la facture de l'investisseur}\label{tut:fm:the-multi-manager-platform:tut}

\begin{tutorial}
\textbf{Objectif.} Chiffrer la structure de frais d'une plateforme et son \omterm{def:fm:the-multi-manager-platform:netting}{risque de compensation des résultats}. \textbf{État final~:}
\cref{fig:fm:the-multi-manager-platform:waterfall,fig:fm:the-multi-manager-platform:netting}.
\begin{tutsteps}
\item \textbf{Pods.} \texttt{firm.podshop.pods(50, 20, 1.0, 0.20, 0.1, rng)} simule les rendements quotidiens des pods~; \texttt{annual} les somme par année.
\item \textbf{Reversements.} \texttt{payouts(ann, 0.20, carry=True, fire=0.10)} rémunère chaque équipe sur son propre profit
(\cref{lst:fm:the-multi-manager-platform:payouts}).
\item \textbf{La facture.} \texttt{fm\_platform.waterfall()} et \texttt{classic()} comparent la structure de \omterm{def:fm:the-multi-manager-platform:pass}{frais refacturés} à un fonds 2-and-20 sur les
mêmes pods.
\item \textbf{Compensation.} \texttt{fm\_platform.netting\_vs\_sr()} compare la simulation à la \cref{prop:fm:the-multi-manager-platform:netting}.
\item \textbf{Échelle et surcouche.} \texttt{ladder\_effect()} et \texttt{overlay\_effect()}.
\end{tutsteps}
\end{tutorial}

\textbf{Pour aller plus loin.} Conservez chaque équipe pour toujours (exercice 7)~; corrélez les pods à 0.4 et observez le gain de la surcouche se réduire dès que seule une partie
du risque commun est un facteur négociable.

\section{Construction~: le modèle de plateforme}\label{bld:fm:the-multi-manager-platform:podshop}

\begin{build}
\bfield{Rôle} L'économie d'une plateforme dans un seul module~: reversements, refacturation des frais, allocation du gérant, \omterm{def:fm:the-multi-manager-platform:netting}{risque de compensation des résultats}, échelle de drawdown et
surcouche du \omterm{def:fm:the-multi-manager-platform:centre}{book central}.
\bfield{Interface} \texttt{firm.podshop}~: \texttt{pods(n, years, sr, vol, rho, rng, days, beta)}, \texttt{annual}~;
\texttt{payouts(ann, rate, carry, fire)}, \texttt{netting\_cost}~; \texttt{expected\_positive(mu, sigma)}~; \texttt{Terms}, \texttt{waterfall(ann, terms, fire)},
\texttt{classic(ann, mgmt, perf)}~; \texttt{ladder(daily, cut, stop)}~; \texttt{overlay(daily, beta, factor)}.
\bfield{Règles} Le reversement d'une équipe ne dépend jamais du résultat d'une autre équipe~; les pertes reportées ne réduisent que ses propres reversements ultérieurs~; le report de carry d'une équipe remplacée est
oublié~; la cascade s'additionne exactement.
\bfield{Tests d'acceptation} \texttt{code/firm/podshop/tests/}~: report de carry et absence de récupération (clawback) sur un exemple à la main~; la forme close contre Monte
Carlo~; moyenne, volatilité et corrélation des pods~; identité de la cascade~; l'échelle coupe puis arrête~; la surcouche réduit la volatilité du fonds~; le remplacement
rétablit les reversements.
\bfield{Pour aller plus loin} Allouer le capital selon le rendement ajusté du risque glissant (\texttt{allocate} de firm.multistrat)~; un reversement garanti la première année pour les nouvelles
équipes~; le \omterm{def:fm:the-multi-manager-platform:centre}{book central} s'ajoutant aux positions consensuelles des pods et son effet sur la concentration.
\end{build}

\begin{omsources}
\item Millennium Management LLC, Form ADV Part 2A brochure, 25 août 2026 (SEC Investment Adviser Public Disclosure).
\item One Quant Book 8, chapitre~28 (pods, drawdown limits), et One Quant Book 7, chapitre~28 (crowding).
\item S.~J.~Grossman et Z.~Zhou, «~Optimal investment strategies for controlling drawdowns~», \emph{Mathematical Finance} 3(3), 1993.
\end{omsources}

\section{Exercices}

\begin{exercise}[$\star$]\label{exo:fm:the-multi-manager-platform:1}
Une équipe gère 500~millions de dollars du capital du fonds et réalise 12~\% de rendement sur un an. Avec un \omterm{def:fm:the-multi-manager-platform:payout}{taux de reversement} de 20~\%, quel montant lui est versé~?
\end{exercise}

\begin{exercise}[$\star$]\label{exo:fm:the-multi-manager-platform:2}
La même équipe a perdu 5~\% l'année précédente et a été conservée. Que reçoit-elle sur les 12~\% de cette année, les pertes étant reportées~?
\end{exercise}

\begin{exercise}[$\star$]\label{exo:fm:the-multi-manager-platform:3}
Le profit annuel d'un pod est gaussien, avec un ratio de Sharpe de 1. Calculer $\E[\max(X,0)]$ en unités de sa volatilité.
\end{exercise}

\begin{exercise}[$\star\star$]\label{exo:fm:the-multi-manager-platform:4}
Dans l'\cref{ex:fm:the-multi-manager-platform:two}, chaque pod coûte 3~millions de dollars par an à faire fonctionner. Calculer ce que l'investisseur paie au total et ce que
le fonds gagne grâce à lui.
\end{exercise}

\begin{exercise}[$\star\star$]\label{exo:fm:the-multi-manager-platform:5}
Utiliser la \cref{prop:fm:the-multi-manager-platform:netting} pour calculer le coût de compensation en part du profit brut, à un \omterm{def:fm:the-multi-manager-platform:payout}{taux de reversement} de 20~\%, pour des pods de
ratio de Sharpe 0.5.
\end{exercise}

\begin{exercise}[$\star\star$]\label{exo:fm:the-multi-manager-platform:6}
L'échelle du chapitre coupe 88~\% des années-équipes. Expliquer pourquoi à partir de la volatilité des pods, et proposer des seuils qui en couperaient environ un quart.
\end{exercise}

\begin{exercise}[$\star\star\star$]\label{exo:fm:the-multi-manager-platform:7}
\emph{Programmation.} Relancer \texttt{fm\_platform.waterfall(fire=None)}~: aucune équipe n'est jamais remplacée. Que deviennent les reversements aux équipes et la part de l'investisseur,
et quel est le coût de compensation en part du profit brut~?
\end{exercise}

\begin{exercise}[$\star\star\star$]\label{exo:fm:the-multi-manager-platform:8}
\emph{Trouver la faille.} «~Le ratio de Sharpe moyen de nos équipes est de 1.5 sur leurs historiques, donc le \omterm{def:fm:the-multi-manager-platform:netting}{risque de compensation des résultats} ne nous coûte presque rien.~»
\end{exercise}

\section{Problème~: payés pour les gagnants}

\begin{problem}[{Problème du week-end --- payés pour les gagnants}]\label{pb:fm:the-multi-manager-platform:1}
Un investisseur se voit proposer une plateforme multigérants comptant cinquante pods et une structure de \omterm{def:fm:the-multi-manager-platform:pass}{frais refacturés}. On dispose de sa plaquette et d'un simulateur.

\medskip\noindent\textbf{Partie I --- La structure.}
\begin{enumerate}
\item Définir des \omterm{def:fm:the-multi-manager-platform:pass}{frais refacturés} et un \omterm{def:fm:the-multi-manager-platform:payout}{taux de reversement}.
\item D'après la plaquette, comment se détermine la rémunération d'un gérant de portefeuille, et que deviennent les pertes d'une équipe~?
\item Que fait le centre du capital du fonds, en dehors de l'allouer aux équipes~?
\item Pourquoi les investisseurs d'une plateforme peuvent-ils verser des primes une année où le fonds perd de l'argent~?
\end{enumerate}

\medskip\noindent\textbf{Partie II --- La facture.}
\begin{enumerate}[resume]
\item Donner la cascade par tranche de 100~\$ de capital~: brut, reversements aux équipes, coûts, gérant, investisseur.
\item Dans combien des 400 années simulées le fonds a-t-il perdu de l'argent~?
\item Que laisserait un fonds 2-and-20 à ses investisseurs sur les mêmes pods, et pourquoi ne peut-il pas exister~?
\item Quelle est la part de l'investisseur dans le profit brut~?
\end{enumerate}

\medskip\noindent\textbf{Partie III --- Compensation.}
\begin{enumerate}[resume]
\item Définir le \omterm{def:fm:the-multi-manager-platform:netting}{risque de compensation des résultats} et énoncer la \cref{prop:fm:the-multi-manager-platform:netting}.
\item Calculer le coût en forme close pour des ratios de Sharpe de 1, 0.5 et 0.25.
\item Donner le coût simulé lorsque les équipes sont remplacées après une perte de 10~\%, pour des ratios de Sharpe de 1 et 0.5.
\item Donner le coût simulé lorsqu'aucune équipe n'est jamais remplacée, et expliquer la différence.
\item Quelles sont les deux propriétés d'une plateforme qui déterminent son risque de compensation~?
\end{enumerate}

\medskip\noindent\textbf{Partie IV --- Le centre.}
\begin{enumerate}[resume]
\item Donner la volatilité et le ratio de Sharpe du fonds avec et sans la \omterm{def:fm:the-multi-manager-platform:centre}{surcouche factorielle}, et expliquer pourquoi le gain est une borne supérieure.
\item Donner la part des années-équipes coupées et arrêtées par l'échelle, ainsi que son effet sur le rendement, la volatilité et le ratio de Sharpe.
\item Pourquoi l'échelle ne coûte-t-elle quelque chose que dans ce modèle~?
\item Définir le \omterm{def:fm:the-multi-manager-platform:centre}{book central} (centre book) et dire qui possède son P\&L.
\item Comment faut-il limiter le risque du \omterm{def:fm:the-multi-manager-platform:centre}{book central}~?
\item Énoncer le \emph{résultat nommé}~: le coût de compensation par pod en forme close, sa valeur pour des ratios de Sharpe de 1 et 0.5, et la part de
pass-through dans le profit brut (reversements et coûts) pour cinquante pods.
\item En deux phrases, que doit demander l'investisseur à la plateforme avant d'investir~?
\end{enumerate}
\end{problem}

\section{Questions d'entretien}

\begin{interviewq}[$\star$ \iqroles{trader, chercheur}]\label{iq:fm:the-multi-manager-platform:1}
Que sont les \omterm{def:fm:the-multi-manager-platform:pass}{frais refacturés}, et pourquoi les plateformes multigérants les utilisent-elles plutôt que des frais de gestion fixes~?
\end{interviewq}

\begin{interviewq}[$\star$ \iqroles{trader}]\label{iq:fm:the-multi-manager-platform:2}
Vous dirigez un pod qui a perdu 5~\% l'an dernier et gagné 12~\% cette année. Comment votre reversement est-il calculé, et que se passe-t-il si vous aviez été remplacé~?
\end{interviewq}

\begin{interviewq}[$\star\star$ \iqroles{chercheur, risque}]\label{iq:fm:the-multi-manager-platform:3}
Deux pods ont le même ratio de Sharpe~: pourquoi rémunérer chacun sur son propre profit coûte-t-il plus cher que de rémunérer sur leur somme~? Donnez la formule.
\end{interviewq}

\begin{interviewq}[$\star\star$ \iqroles{risque}]\label{iq:fm:the-multi-manager-platform:4}
Cinquante pods de volatilité 20~\% chacun, avec une corrélation deux à deux de 0.1~: quelle est la volatilité du fonds, et qu'en serait-il si la part commune était couverte~?
\end{interviewq}

\begin{interviewq}[$\star\star$ \iqroles{trader}]\label{iq:fm:the-multi-manager-platform:5}
Pourquoi un arrêt à un drawdown de 10~\% est-il une règle sévère pour un pod de volatilité 20~\%~? Que proposeriez-vous~?
\end{interviewq}

\begin{interviewq}[$\star\star\star$ \iqroles{chercheur}]\label{iq:fm:the-multi-manager-platform:6}
Une plateforme affirme que son \omterm{def:fm:the-multi-manager-platform:centre}{book central} ne fait que couvrir. Qu'examineriez-vous dans son P\&L et ses positions pour tester cette affirmation~?
\end{interviewq}
